\chapter{Industrialisation de la plateforme}
\label{ch:industrialisation}

Une fois la chaîne fonctionnelle, plusieurs limites ont été identifiées : absence d'historique,
reconstruction complète de toutes les tables à chaque exécution, application dépendante d'un poste
local, absence d'alerte en cas d'incident et validation manuelle des modifications. Ce chapitre présente
les réponses apportées. L'historisation ayant été décrite à la section~\ref{sec:scd2}, il traite les
quatre autres points.

\section{Traitements incrémentaux}
\label{sec:incremental}

Dans la version initiale, \texttt{dbt build} reconstruisait intégralement toutes les tables. Ce coût,
négligeable sur WWI, devient prohibitif sur des volumes réels. Les \textbf{dix modèles les plus
volumineux} sont désormais matérialisés en mode \texttt{incremental} avec la stratégie \texttt{merge}
(tableau~\ref{tab:incremental}).

\begin{table}[H]
\centering
\caption{Modèles incrémentaux}
\label{tab:incremental}
\begin{tabular}{@{}l l l@{}}
\toprule
\textbf{Couche} & \textbf{Modèle} & \textbf{Clé de fusion} \\
\midrule
\multirow{6}{*}{Silver} & \texttt{slv\_invoice\_lines} & \texttt{invoice\_line\_id} \\
 & \texttt{slv\_invoices} & \texttt{invoice\_id} \\
 & \texttt{slv\_order\_lines} & \texttt{order\_line\_id} \\
 & \texttt{slv\_orders} & \texttt{order\_id} \\
 & \texttt{slv\_stock\_item\_transactions} & \texttt{stock\_item\_transaction\_id} \\
 & \texttt{slv\_customer\_transactions} & \texttt{customer\_transaction\_id} \\
\midrule
\multirow{4}{*}{Gold} & \texttt{fact\_sales} & \texttt{invoice\_line\_id} \\
 & \texttt{fact\_orders} & \texttt{order\_line\_id} \\
 & \texttt{fact\_stock\_movements} & \texttt{stock\_item\_transaction\_id} \\
 & \texttt{fact\_customer\_transactions} & \texttt{customer\_transaction\_id} \\
\bottomrule
\end{tabular}
\end{table}

Une macro \texttt{incremental\_filter} ne retient que les lignes dont la date de dernière modification
est postérieure à la plus récente déjà présente dans la table cible. Pour les faits issus d'une jointure
en-tête/ligne, la date retenue est la plus récente des deux (\texttt{GREATEST}), de sorte qu'une
modification de l'en-tête (par exemple la date de livraison d'une facture) est aussi propagée.

\begin{resultbox}[title=Vérification]
Après une reconstruction complète (\texttt{dbt build -{}-full-refresh}), une exécution normale des dix
modèles incrémentaux a traité \textbf{0 ligne} au lieu de 228\,265 pour \texttt{fact\_sales} : seules les
modifications seront désormais retraitées.
\end{resultbox}

\section{Alertes et supervision}

\subsection{Alertes Dagster}
Le module \texttt{orchestration/alerts.py} ajoute deux capteurs actifs par défaut :
\begin{itemize}
  \item \texttt{wwi\_run\_failure\_alert} : un \emph{run failure sensor} déclenché dès qu'un run échoue.
        Il construit un message (job, identifiant du run, étapes en échec, erreur, lien vers
        l'interface) et l'envoie sur \textbf{Slack} (webhook) et/ou par \textbf{e-mail} (SMTP) selon
        la configuration du fichier \texttt{.env} ; à défaut, l'alerte est journalisée ;
  \item \texttt{gold\_freshness\_sensor} : des \textbf{contrôles de fraîcheur} vérifient que les tables
        Gold principales ont été rafraîchies au cours des 26 dernières heures (exécution quotidienne
        plus une marge de deux heures). Un planning bloqué ou un démon arrêté est ainsi détecté.
\end{itemize}

\subsection{Maîtrise des coûts Snowflake}
Un \textbf{moniteur de ressources} (\texttt{WWI\_CREDIT\_MONITOR}) est associé à l'entrepôt
\texttt{WWI\_WH} : il envoie une notification à 75\,\% et 90\,\% du quota mensuel de crédits, suspend
l'entrepôt à 100\,\% et l'arrête immédiatement à 110\,\%. Il protège en particulier les crédits limités
du compte d'essai.

\section{Déploiement dans Streamlit in Snowflake}

L'application, initialement exécutée sur le poste du développeur, a été rendue déployable dans
\textbf{Streamlit in Snowflake}. Ce mode d'hébergement présente plusieurs avantages :
\begin{itemize}
  \item aucun serveur ni identifiant à gérer : l'application utilise la session de l'utilisateur
        connecté à Snowsight ;
  \item les données ne quittent pas Snowflake ;
  \item le partage se fait par les mécanismes de droits de Snowflake.
\end{itemize}

Le script \texttt{deploy/deploy\_sis.py} automatise le déploiement :
\begin{enumerate}
  \item création du schéma \texttt{WWI\_DB.APP} et attribution au rôle \texttt{WWI\_ANALYST} des droits
        de création de stage et d'application ;
  \item création d'un stage interne et envoi des fichiers de l'application (\texttt{PUT}) : code,
        modèle sémantique, logo et fichier \texttt{environment.yml} des dépendances ;
  \item création de l'objet \texttt{STREAMLIT WWI\_DB.APP.WWI\_INSIGHTS} avec l'entrepôt
        \texttt{WWI\_WH}.
\end{enumerate}

L'application appartient au rôle \texttt{WWI\_ANALYST} : elle s'exécute donc avec ses seuls droits de
lecture sur Gold. Le même code source fonctionne dans les deux environnements, le moteur détectant la
présence du module interne \texttt{\_snowflake} pour choisir entre l'appel REST local et l'API interne
\texttt{\_snowflake.send\_snow\_api\_request}.

\section{Intégration et déploiement continus}

Le dépôt est hébergé sur GitHub. Le workflow \texttt{.github/workflows/ci.yml} s'exécute à chaque
\emph{push} et à chaque \emph{pull request} (figure~\ref{fig:ci}).

\begin{figure}[H]
\centering
\resizebox{\linewidth}{!}{%
\begin{tikzpicture}[node distance=1.1cm]
  \node[box, fill=gray!10] (push) {\textbf{push / PR}\\ GitHub};
  \node[box, fill=blue!8, right=of push, text width=3.3cm] (q) {\textbf{Qualité}\\ \scriptsize ruff, compilation,\\ validation YAML};
  \node[box, fill=snow!15, right=of q, text width=3.5cm] (d) {\textbf{dbt build}\\ \scriptsize sur clone zéro copie\\ \texttt{WWI\_DB\_CI\_<run>}};
  \node[box, fill=green!10, right=of d, text width=3.3cm] (dep) {\textbf{Déploiement}\\ \scriptsize Streamlit in Snowflake\\ (branche \texttt{main})};
  \draw[flow] (push) -- (q);
  \draw[flow] (q) -- (d);
  \draw[flow] (d) -- (dep);
  \node[below=0.35cm of d, font=\scriptsize, text width=4.5cm, align=center] {le clone est supprimé à la fin,\\ même en cas d'échec};
\end{tikzpicture}}
\caption{Chaîne d'intégration et de déploiement continus}
\label{fig:ci}
\end{figure}

\begin{enumerate}
  \item \textbf{Qualité} : analyse statique du code Python (\texttt{ruff}), compilation de tous les
        scripts, validation syntaxique de tous les fichiers YAML et de la structure du modèle sémantique.
  \item \textbf{dbt} : la base de production est clonée par \textbf{clonage zéro copie}
        (\texttt{CREATE DATABASE WWI\_DB\_CI\_<run> CLONE WWI\_DB}) : le clone est instantané et ne
        duplique aucune donnée. \texttt{dbt build} y est exécuté (modèles, snapshot et tests), puis le
        clone est supprimé. Les modifications sont ainsi testées sur des données réelles sans jamais
        toucher à la production. La base cible étant paramétrée par la variable
        \texttt{SNOWFLAKE\_DATABASE}, aucun code n'est spécifique à la CI.
  \item \textbf{Déploiement} : sur la branche \texttt{main}, si les étapes précédentes réussissent,
        l'application Streamlit in Snowflake est redéployée.
\end{enumerate}

Les identifiants Snowflake sont stockés dans les \emph{secrets} chiffrés du dépôt GitHub, jamais dans
le code.
